\hypertarget{homejimdocumentsmedicalkidney_glossary.md}{%
\paragraph{/home/jim/documents/medical/kidney\_glossary.md}\label{homejimdocumentsmedicalkidney_glossary.md}}

\hypertarget{as-of-7-sept-2020}{%
\subsection{As of 7 SEPT 2020}\label{as-of-7-sept-2020}}

\hypertarget{section}{%
\paragraph{}\label{section}}

NIH glossary; AUA 2014 Guidelines, www.kidney.org * ANATOMY of KIDNEY-
http://www.healthpages.org/anatomy-function/kidney/ *
\href{https://www.cdc.gov/diabetes/programs/initiatives/kidney.html}{CDC}
(lot here) * U of Chicago (Coe) http://kidneystones.uchicago.edu/ * OHSU
text: ``Nephrology Secrets'' \textbar{} ``Manual of Nephrology'', Schier
\textbar{} WJ300 .M294 * ``Renal and \ldots{} Disorders'', 8th, Schier
\textbar{} WJ300, R391,1018 * Netter: Clinical Anatomy \textbar{} QS 17
* Chapple, et al ``Urodynamics Made Easy \ldots{}'' \textbar{}
WJ102.C467 u * Smith's \textbar{} WJ100 .S64 2013 * Medscape: Nice
overview of Hypocitraturia (http://bit.ly/2iLE7HD)

Heading Level 3, for letters. Normal text is Arial 9. For Single
spacing, Shift-Enter.

Papers Asplin, J.R. Urolithiasis (2016) 44: 33.
https://doi.org/10.1007/s00240-015-0846-5 (Litholink) he management of
patients with enteric hyperoxaluria.

-- A --

ACIDOSIS (vs.~ALKALOSIS) - Acidic blood. Normal Blood is basic pH-
7.35-7.45 (Urine s/d be slightly acidic, 6.5-7.0)
https://en.wikipedia.org/wiki/Acidosis Use KCit, NaCit; Na Bicarbonate
to lower blood acid. Bicarbonate (HC03 -1 is blood pH buffer. Can add H
or lose H) .

\begin{itemize}
\item
  Metabolic Acidosis - Kidney not removing all acids from blood, or not
  enough bicarbonate produced to alkalize.. Blood pH drops. Blood starts
  to absorb Citrate (to alkalize) from Kidney, leaving less Cit for
  urine. Metabolic.
\item
  Acidosis - acidic blood- . is bad for brain/heart.
\end{itemize}

Acidosis (AUA) can arise from a diet that is inordinately rich in foods
with a high potential renal acid load compared to low-acid (i.e.,
alkaline) foods.83 (high protein diet, forces kidney adaptation to
handle acid load). Foods providing an acid load include meats, fish,
poultry, cheese, eggs, and to a lesser extent, grains. Foods conferring
an alkali load include nearly all fruits and vegetables. Milk and
yogurt, as well as fats, are essentially neutral for acid load.84 KCit
therapy attempt to return kidney to normal state, freeing Citrate to
flow to urine.

Acid Load - also increases UrCa excretion (depletes bones), esp sulfur
amino acids (mostly animal), such Methionine, Cysteine

\textbf{ALKALI THERAPY} - Use of one of : K+Citrate (not always
tolerated, 1 mEq = 1080 g), Mg+Citrate (Magnesium), Na+Citrate (usu well
tolerated but sodium). Goal is increase urinary citrate excretion
(\textgreater{} 3 mmol/day).

\begin{itemize}
\tightlist
\item
  Dose large, up to 120 mEq/day (Me: 40 mEq= 4 * 10mEq per pill).
\item
  Caution: Western diet (protein) acid load to kidney. Kidney adjusts.
  With Alkali therapy, acid-base balance roughly returns. Kidney can
  give up its modifications to acid load. BUT, balance not exact.\\
\item
  See Coe: How Citrate gets into Urine.
\end{itemize}

ALLOPURINOL - -- Dosing issues,
https://www.ncbi.nlm.nih.gov/pubmed/17897242 -- From Kaiser:
``Allopurinol is used to treat gout and certain types of kidney stones.
It is also used to prevent increased uric acid levels in patients
receiving cancer chemotherapy. These patients can have increased uric
acid levels due to release of uric acid from the dying cancer cells.
Allopurinol works by reducing the amount of uric acid made by the body.
Increased uric acid levels can cause gout and kidney problems.''

ANION GAP -- (mine=?) Usually in blood, but can be urine or any fluid.
Measures the excess of specific cations (Na+, K+ etc) over specific
anions (Cl-, HOCK. K+ is lower and often ignored (?) HOCK- is
bicarbonate or ``CO2''. (details) Why positive, not zero? AG measures
presence of OTHER (missing) cations, an indicator of disease?

APRT Deficiency (rare) -- Adenine (nucleobase) normally converts to AMP.
But APRT is genetic error, produces wrong enzyme, adenine converts to
wrong molecule which is toxic to kidney. May not appear until adult.
15-20\% do not know they have it. Effect varies even within same family.
(https://ghr.nlm.nih.gov/condition/adenine-phosphoribosyltransferase-deficiency)

-- B --

\textbf{BILE ACID MALABSORPTION} --
https://en.wikipedia.org/wiki/Bile\_acid\_malabsorption Effect can
diarrhea. Not easy to test.

\textbf{B12} - not available in plant form; best sources: fish (clam,
tuna, trout), meat (beef liver), dairy also. Role?

BHP -- Beign \ldots. Via UROFLOW + IVP can yield lots of info? BLADDER
-- volume \textasciitilde{} 400 ml (me: pee \textasciitilde{} 100-200
ml?) - Chronic Residual Urine (?) - reduced bladder function? - Acute
Cystitis - Bladder can not stretch , involuntary pee see Smith's BLOOD -
SERUM - BUFFER - maintains pH by binding/releasing proton as needed.
Bicarbionate (HC03 -1) Favorite buffer (for humans). -- BiC- + H+
-\textgreater{} Carbonic Acid → CO2 (and reverse).

\textbf{BUN} (me = ?)-- protein metabolism. Imperfect measure of
filtration rate because affected by GFR, protein intake, hydration,
gastric bleeding and can be fooled with too much water. - Normal BUN/Cr
= 10/1; obstruction : 20-40/1.

BUTYRATE - short-chain fatty acid, produced by gut bacteria CLAIM: in
large intestine, tightens protein chains that bind cells together,
attracts mucus coating (good), inhibits `leaky gut' (thought disproven)
.

-- C --

\textbf{CaOx - 2 crystal forms:} CaOx mono (COM) and CaOx dihydrate (CO
), the later worse because binds to epithelial cells. Proteins act as
`glue' to hold CaOx monohydrate (COM) crystals, once formed, so COM
aggregates into a stone.

(True?) -- ME: 2JAN2014: 50\% Mono; 40\% Dihydrate, 10\% calcium
phosphate (hydroxy- and carbonate-) \textbar{} 41 mg \textbar{} 5-9 mm.
* Stone Analysis: Claim: If mostly monohydrate, suspect PH. * If mixed,
suspect SH.

CHELATING - an agent (Ex: citrate) inhibits Ca crystal formation in
kidney. 2 Cit bind 3 Ca(+2) atoms. Several proteins also inhibit - at
least 6 times, not well understood.

\textbf{CHRONIC KIDNEY DISEASE, CKD} (vs.~ACUTE) - Ongoning, gradual
loss of kidney function. (many causes - from diabetes to xxx prolonged
urinary obstruction) DIET: Less protein (esp non-dairy animal), less Na,
less K. Good news - loss can be controlled with proper care. About 15\%
US with CKD. Most die WITH CKD, not from. * Stage I, GFR\textgreater{}
90, but kidney damage. * Stage II, GFR 60-89, but kidney damage. * Stage
III, GFR less than 60 for 3 months, even if no sign of kidney damage
(protein, mcroalbum\ldots).\\
* Stage IV, GFR less than 30? Nephrologist becomes primary doctor . *
Stage V, GFR less than 15. Dialysis or Transplant necessary for life.

\begin{itemize}
\item
  Causes: Diabetes (38\%) \textbar{} Hypertension (24\%) \textbar{}
  Glomerular Disease (15\%) \textbar{} Polycystic (5\%) \textbar{} Other
  (18\%) -- Symptoms of decreasing KIdney Function:
\item
  Acidosis (acid in blood, should be base) EDEMA - swelling of tissues,
  esp
\item
  ankles, due to fluid (Na) build up; can occur BEFORE proteinuria
  appears?
\item
  ERYTHROPOIETIN - too few new RED blood cells ⇒ anemia Phosphate build
  up.
\item
  Kidneys unable to remove all and Ca leaches from bones; Ca Phosphate
  forms.
\end{itemize}

(ToDo: rewrite) \textbf{CITRATE} (wiki:
https://en.m.wikipedia.org/wiki/Citric\_acid) (see Coe: urine and
citrate)

Citric acid is weak acid, stronger than acetic. pH is a measure of
equilibrium H+ concentration. However, most common, ie weak acids, not
strong enough to be fully dissociated, so H+ not quite as high as
expected, giving higher PH.

Citrate dissociates in 3 possible polyatomic forms: - Citrate anion (-3)
is form with ability to sequester Ca. (2 Citrate - molecules bind 3 Ca
in a SALT) Conc of this form varies with pH.

Citric Acid itself does NOT bind with Ca+2 and by itself is NOT useful
to reduce Ur Ox.

\begin{itemize}
\item
  Citrate vs Citric Acid - pH dependent! See Why Lemon/Lime Juice in the
  form of Citrate (-3), while Orange juice is more Citric Acid form, pH
  6+ Citrate dominates
\item
  https://www.physicsforums.com/threads/concentation-of-citric-acid-in-fruit-juice.428241/
\item
  https://www.physicsforums.com/threads/ph-citric-acid.164118/
\end{itemize}

\textbf{Urine Citrate} - Regulated by Kidney. To prevent Ca stones, risk
is urine citrate dependent (not urine pH) Higher citrate, lower calcium
in urine, lower the risk. ( But uric acid stones, the risk is pH
dependent: goal is \textasciitilde pH=6.) KIDNEY tries to balance URINE
ACIDITY (explain) (normal urine pH=5.5-7.0, better is 6.5-7.0)

Kidney also produces or absorb citrate. Per AUA, healthy urinary citrate
\textgreater{} 640 mg/day (population mean) Debate: Is it Citrate per
day, or concentration? pH role? To know Citrate content in a fluid, ,
need to know pH and total amount of citric acid.

Do Not Believe diet soda citrate Label (Coe). Diet 7-up (10 mEq/L) and
Crystal Light Lemonade (21.7 mEq/L)\ldots{} are good sources of CITRATE!
1 KCitrate pill=10mEq

CaCitrate exists in 3 crystal forms (1 is rare) Avoid DIET drinks to
supplement Citrate. Crystal Light Lemonade has 21.7 mEq Citrate/L (=2
KCitrate pills), Diet 7-Up = 10 mEq/L Reason: Per NURSES STUDY, above
age 40 GFR decreases \textasciitilde{} 1 mL/min per year, normal. 1 such
DIET drink per day is no problem. But those who drank 2 such DIET drinks
per day, lost on average 3 mL/min per year.

Tamarind is excellent Citrate source, Tartaric Acid= sour taste,\\
HCA = Citric + extra OH? (derivative of citric acid) , Citrate Lyase
(https://en.wikipedia.org/wiki/Tartaric\_acid) (see:
http://www.news.pitt.edu/news/researchers-propose-new-treatment-prevent-kidney-stones)

Lemon Juice in acidic solution is primarily citric acid, not in Citrate
form (Coe) At higher pH( 6.5), Citrate form dominates Vitamin C,
ascorbic acid, has nothing to do with Citric Acid, though both are found
in citrus fruits. In blood, form of citric acid is CITRATE, 80-170
umol/L. (say \textasciitilde120 umol/L = 0.12 mmol/L) 5-20\% will wind
up in urine. Kidney will reabsorb rest. In blood, also have free Ca+2
(cations, and not bound to protein), \textasciitilde{} 1.0 mmol/L. Also
free Mg. Both compete to bound with Oxalate. CITRATE is stable (here is
where PTH comes in?) Blood pH (7.35-7.45) is well-regulated, homeostis.
In body, CITRATE affects LIPIDS, BRAIN, GLUCOSE, transport into CELLS.
NaDC1? NaDC2? Reference
https://www.physicsforums.com/threads/ph-citric-acid.164118/

CITRATURIA - presence of citrate in urine. I have HYPOCITRATURIA
(def:less than 320 mg/day ). Stone-formers are told to keep 500-800
mg/day.

\begin{itemize}
\tightlist
\item
  CLAIM: purpose of kidney Citrate production is balance urine acidity.
  If low urine pH (acid), citrate levels are maintained and kidney
  produces less. If high urine pH (base), kidney will increase citrate
  production. Urine pH usually slightly acidic (5.5-7.0; desired is
  6.5-7.0) https://en.wikipedia.org/wiki/Urine\#pH
\end{itemize}

COLLAGEN -- protein. Tendons, bone-muscle, and Ligaments, bone-bone, are
collagen.

\textbf{COLESTYRAMINE} (https://en.wikipedia.org/wiki/Colestyramine) -
Colestyramine removes bile acids from the body by forming insoluble
complexes with bile acids in the intestine, which are then excreted in
the feces. Normally bile acides reabsorbed near terminal portion of the
small bowel \textbf{(ileum)}

**CREATININE - Waste product in blood, kidney should remove. Urinary
Creatinine excretion is relatively stable, day-to-day.
\textasciitilde8-20g/kg. Product of muscle breakdown. Kidneys remove
almost all creatinine from blood, without any modification. Rise in
blood creatinine \emph{may} be signal of reduced kidney function (or ate
more protein, intense exercise etc than usual).

-- D --

DIGESTION -- 4-6 hours after eating -- stomach empties Additional 6-8
hours (ie 10-16 hours after eating) -- intestine processes, dumps into
colon 1-3 days in Colon

TEST: eat corn. When do indigestible kernels appear in stool? Moral? If
you get sick, probably what you ate a day or so ago, not the most
recent. Diuretic - parsley!

-- E --

\textbf{EDEMA (``dropsy'')} - fluid swelling, typically around ankles.
(Many possible causes and many other uses for EDEMA.) In case of kidney
issues, EDEMA will arise before protein appears in urine (proteinuria)
and is due to compromised kidney glomeruli or its capillaries??

\textbf{EGD} (no!(mistake) I will have 14 July 2017, with colonoscopy)
wiki: https://is.gd/WyWgU6 Can rule out many things, but does not
diagnose IBS Medline: https://is.gd/VkQFSA Inspection of stomach to 2nd
part of 4 of duodenum, which itself is 1st of 3 parts to small
intestine. (wiki: https://is.gd/kCOnyD) Normal: The esophagus, stomach,
and duodenum should be smooth and of normal color. Sounds like looking
for irritation, ulcers, anything other than smooth tissue. This may
include one of several issues from Crohn's disease (an autoimmune
response, one of Inflamatory Bowel Diseases) to \ldots. Symptoms of IBD
(which I don't have) can indicate one of several problems, which visual
inspection can detrmine (no blood test, for example).

ENDOCRINOLOGY - (wiki) study of glands and their secretions,
particularly the \emph{diseases} related to ETIOLOGY - underlying cause

\textbf{ENTERIC} - relating to intestines; passing through stomach to
intestines; ex: enteric aspirin.

\textbf{ENTERIC HYPEROXALURIA} Fat malabsorption promotes free fatty
acids to bind with diet calcium, thus reducing calcium available to
precipitate diet oxalate. Delivery of unabsorbed bile salts and fatty
acids to the colon increases colonic permeability, the site of oxalate
hyper-absorption in enteric hyperoxaluria. The combination of soluble
oxalate in the intestinal lumen and increased permeability of the
colonic mucosa leads to hyperoxaluria.

ENDOGENOUS Ox - Ox origin is inside body, e.g.~synthesis, liver is main
organ

Endourolologist - Urologist concentrating in urolithiasis, including
trans-ureter procedures (vs laparoscopy?)

\textbf{ERYTHROPOIETIN} - too few new RED blood cells ⇒ anemia EXOGENOUS
- Ox origin is outside, e.g.~food; Vit C breakdown is considered a 3rd
source of Ox

-- F --

\textbf{FAT MALABSORPTION }- - Normal: fat is absorbed in small
intestine, little reaches COLON.

\begin{itemize}
\item
  Abnormal: Fat reaches COLON. Since FAT tends to combine with Ca
  (why?), Ox remains free risking, absorption in colon. Many possible
  reasons. Hard water (https://en.wikipedia.org/wiki/Hard\_water) is
  rich in minerales like Mg, Ca. Residue forms with Ca(`goo') when using
  hand soap (fatty acid). Note: Small intestine removes all from
  \textasciitilde{} 0.5 g/d, however normal feces up to 7 g fat. see:
  see https://en.wikipedia.org/wiki/Fecal\_fat\_test.
\item
  Normally up to 7 grams of fat can be malabsorbed in people consuming
  100 grams of fat per day. In patients with diarrhea, up to 12 grams of
  fat may be malabsorbed since the presence of diarrhea interferes with
  fat absorption, even when the diarrhea is not due to fat
  malabsorption.
\end{itemize}

see: https://en.wikipedia.org/wiki/Saponification - Related (?) Fatty
Acid: Arachidonic Acid (20 chain) traps Ca+, Mg+ in intestine, leading
to increased UrOx

FIBER (or PLANT STRUCTURE?) Acts as LAXATIVE to speed up excretion
(decrease transit time). Unavailable to digestion via human intestinal
enzymes. Source: plant cell walls (97-98\% of diet) or polysaccharides
(2-3\%, food additives as pectin, gum) Polymers pass small intestine to
colon. Fermentable or non-Fermentable. Fiber usuallly complex
polysaccharides (lignin proteins) or non-carbohydrate esters, even Ca+2
(?) UK diet \textasciitilde{} 20g/day mostly cerals. African diet
\textasciitilde60-150g/day!

\begin{itemize}
\tightlist
\item
  SEEDS
\item
  BRAN - outer layer - mostly cellulose -indigestible fiber; also
  minerals, vit, protein
\item
  GERM - embryo - rich in unsaturated fats, hence not very stable and
  usually removed in modern milling (roller pressure can melt fats and
  remove)
\item
  ENDOSPERM - starch cells
\item
  RESISTIVE STARCH (wiki) Resistant to amylolytic degradation in
  cooked/processed foods. Not digested by human gut; increasingly found
  useful to feed microbiome which produce host of beneficial effects.
  African diet (beans, grains) 3-4 x Chinese or Indian diet, which in
  turn are multiples of western diets. Bacteria extract energy and
  nutrition from these carbohydrates via fermentation.
\item
  STARCH - glucose storage, branched molecule (amylose) or unbranched
  (amylopectin)
\item
  LIGNAN - wood, bran, cereals - with structure - do not rot easily
\end{itemize}

FODMAPS -- (wiki: https://en.wikipedia.org/wiki/FODMAP) Fermentable
Olgasacchardies, Disaccharides , Monosaccharides, and polygols. If
excess intake, small intestine unable to process and wind up in colon
causing trouble (?) Many healthy foods have; broccoli brussel sprouts
raddico, avacado, peach , mushroom, whole grain, legumes (mimicks
Irritable Bowel Syndrome?) FREE RADICAL - very reactive; can introduce
unintended reactions with DNA and damage. However, some forms are
necessary for life and can kill bacteria.

(TODO: Begin here) FRUCTOSE - ½ glucose, ½ sucrose; only carb that can
increase Ur Ox, Ur Uric Acid, Ur Ca. b/c only carb processed by LIVER,
which controls Ox production. (Be sure to get enough B6, Mg) Main danger
is high fructose syrup HFCS, but any risk from too many apples or plain
yogurt ? In units of per 100 g, Careful with apples \& pears (raisins,
oranges?) (see: 2008 Taylor, Fructose Consumption) apples have 10.4 g
sugar (5.9 free fructose, 2.4 free glucose, 2.1 sucrose) Pears 9.8 g
sugar (6.2, 2.8, 0.8) Onion 5.0 g (2.0, 2.3, 0.7) Sweet Potato 4.2 (0.7,
1.0, 0.7) Carrot 4.7 (0.6,1.0,0.7) Honey 100.0 (50, 44, 1) HFCS 100.0
(50,50,nil) -- and UrCalcium spike, see Coe: http://bit.ly/2ixgvGt
Claims: Glucose bad, but Fructose is worst. UrCal spikes AND UrVol
decreases (double whammy). During the spike, kidney can reabsorb Calcium
as usual. Forces bone to release calcium to maintain balance. No more
than 10\% added sugar. In male, at 2000 cal, this is 50 g added sugar
per day (4 cal/g)

-- G --

\textbf{GFR} - See calculator:
https://www.kidney.org/professionals/KDOQI/gfr\_calculator. Exist models
for plasma oxalate as function of eGFR, etc. (regressions)

GLOMERULUS - filter, entry to kidney tubular filtering system. Sensitive
to blood pressure, scarring, inflammation; supposed to block larger
molecules (blood cells, proteins) while smaller molecules pass.
https://en.wikipedia.org/wiki/Glomerulus\_(kidney) Large surface
area/thickness? Once past Glomerulus, network of tubules, receptors
re-capture anything useful (nutrients, salts). Filtered blood returns
via Renal Vein, while waste collects as urine. Unfiltered blood arrives
via Renal Artery.

\begin{verbatim}
    Water -                         0.28 nm
    DNA-diameter -                 2.5 nm
    Hemoglobin -                 7.0 nm
    E.Coli        -                 2000 nm
    Liver Cell -                 20,000 nm
\end{verbatim}

GLYOXYLATE - precursor to OX, by adding Ca+ Amino acid HYDROXYPROLINE,
very abundant in meat diet (collagen) , catabolizes to Glyoxylate and
pyruvate.

GOUT - uric acid excess/crystallization in blood. (FeOx insoluble) Esp
in BEER, rich in Fe \& Ox.\\
GUT - Long path from mouth to anus.

-- H -- HEMOGLOBIN - normal 14.0 g - 18.0 g/dL (??) HEMATOCRIT - percent
of blood volume due to Red Blood Cells (normal: 40-54\%)

HYDROXYPROLINE protein, found in collagen, and then normally broken down
by liver chemicals. If not, precursor to oxalate, PH. Meat diet is high
in Hydroxyproline. https://clinicaltrials.gov/ct2/show/NCT0203854 (Mayo)

``The purpose of this study is to determine the contribution of
hydroxyproline metabolism to urinary oxalate and glycolate excretion in
patients with primary hyperoxaluria. Oxalic acid (COOH)2 is an end
product of metabolism that is synthesized mainly in the liver. We have
estimated that 10 - 20 mg is synthesized in the body of healthy adults
each day (1). The main precursor of oxalate is glyoxylate (CHO•COOH).
The bulk of the glyoxylate formed is normally transaminated to glycine
(NH2•CH2•COOH) by alanine: lyoxylate aminotransferase (AGT) or reduced
to glycolate (CHOH•COOH) by glyoxylate reductase (GR). Less than 10\% of
the glyoxylate is oxidized to oxalate by lactate dehydrogenase (LDH). In
individuals with the disease, primary hyperoxaluria, AGT, GR, or HOGA
enzyme is deficient and the amount of oxalate synthesized by the liver
increases to 80 - 300 mg per day. The increased oxalate excreted in
urine can cause damage to kidney tissue. Calcium oxalate stones may form
in the kidney or calcium oxalate crystals may deposit in renal tubules
and the renal parenchyma (nephrocalcinosis). An increased rate of
oxalate synthesis could also contribute to idiopathic calcium oxalate
stone disease. Understanding the pathways of endogenous oxalate
synthesis and identifying strategies that decrease oxalate production
could be beneficial for individuals with these diseases. Hydroxyproline
is the primary source of glyoxylate identified in the body (2). Daily
collagen turnover of bone results in the formation of 300 - 450 mg of
hydroxyproline, which cannot be re-utilized by the body and is broken
down. This metabolism yields 180 - 250 mg of glyoxylate. Further
hydroxyproline is obtained from the diet, primarily from meat and
gelatin-containing products. The bulk of the glyoxylate formed is
converted to glycine by the liver enzyme AGT, some to glycolate and a
small amount to oxalate. The proportion of these metabolites is not
known with any certainty. In this study, a quantitative estimate of the
metabolites formed will provide estimates of the contribution of
hydroxyproline turnover to daily oxalate production. These experiments
will provide valuable information for the future assessment of the
contribution of hydroxyproline metabolism to oxalate production in
individuals with primary hyperoxaluria.''

HERNIA -Abdominal may be from past surgery, a sac forms between lining
of abdominal (aka Incisional or Ventral)

REF:
https://www.sages.org/publications/patient-information/patient-information-for-laparoscopic-ventral-hernia-repair-from-sages/

HYPERCALCEMIA - High Serum Calcium. Could be related to PTH glands or
other factors, including excess D3 and Ca supplementation.(May 2020, I
reduced D3 from \textgreater25k IU/week to 15k IU/wk and Nov Litholink
reduced serum Ca.) REF:
https://www.mayoclinic.org/diseases-conditions/hypercalcemia/symptoms-causes/syc-20355523

HYPERCALCIURIA - (2nd most common) Normal Ur Ca
\texttt{\textless{}\ 200\ mg/d\ -\/-\ moderate:\ \textgreater{}200};
severe \textgreater{} 275. (me: 177, 206, 296-2019) Causes: Genetic,
excess Ca absorption intestine or kidney overload; Renal, ``leak of Ca?
Resorptive. Coe: primary hyperparathyroidism (PHPT) and kidney unable to
remove Ca. Take Thiazide. - Secondary Hypercalciuria, could be gland or
?? (see Coe)

HYPOCITRATURIA - less than 320 mg Ur citrate/d (Normal: 290-710). Claim:
Ur Citrate decreases stone risk by binding to Ca first (or inhibit Ca-Ox
bind). Citrate useful to carry metals (Na, K, Ca). LIVER usually absorbs
Citrate (and little reaches Ur?) 300 mg Citrate x 1 mole/192.19 g = 1.5
mmol.\\
Good overview = - Medscape (http://bit.ly/2iLE7HD) (owned by WebMD)
AUA:Urinary citrate is a potent inhibitor of calcium stone formation.80
Although hypocitraturia is variably defined, most healthy individuals
excrete at least 600 mg daily in the urine, and many believe this
urinary level should constitute the minimum target output for stone
formers. Hypocitraturia is a common risk factor for stone disease with
an estimated prevalence of 20-60\%.81, 82 -- Est. 40-60\% stone formers
have this condition.\\
- Treatment for Idiopathic HypoCitraturia (me) = mild metabolic
alkalosis (in urine?) so kidney increases UrCitrate, But keep urine pH
6.5-7.0. 80-90\% successfully raise UrCit this way. lemons, oranges,
grapefruit, and lime (citrus) and melon (non-citrus) good sources of
citrate that will raise Ur\_Citrate.\\
-- Decrease UrCit: strenuous exercise, water pill, animal diet (b/c
lowers blood pH, kidney must retain more Cit)

\textbf{HYPEROXALURIA} (found \textasciitilde20\%)- excess Ur Ox. Normal
\textasciitilde40 mg/d.~Est 60\% Ox by liver; 40\% from diet (varies).
Note: several papers mention that if UrOx is high, look for Primary or
Enteric factors even if not so obvious. (check: normal urinary Ox :
urinary calcium \textasciitilde{} 1:10 - true?) ! ``Most patients with
severe secondary hyperoxaluria have gastrointestinal diseases that cause
fat malabsorption'' {[}Mayo{[}1{]}{]}

\begin{enumerate}
\def\labelenumi{\arabic{enumi}.}
\tightlist
\item
  Primary(1/85,000) (\textgreater{} 100 mg/d, \textgreater1.0
  mmol/1.73m3/d) genetic (rare) - take B6 (pyridoxine) {[}`extremely
  high' - 87-231 mg/d{]}

  \begin{enumerate}
  \def\labelenumii{\arabic{enumii}.}
  \item
    PH1 - test URINE GLYCOLATE (UrOx \textgreater{} 0.8 mmol/24 hours =
    70 mg/24 hrs) (affects young)
  \item
    PH2- test URINE GLYCERATE (affects later in life than PH1) --
    deficiency of the enzyme glyoxylate reductase/hydroxypyruvate
    reductase (GR/HPR) (paper or paper)
  \item
    PH3 - ?
  \item
    Determination: Urine analysis, here; Diagnostic path here
  \end{enumerate}
\end{enumerate}

\begin{itemize}
\tightlist
\item
  (not due to diet); rare metabolic defect, inherited, appears young; 3
  or more types? I don't have BUT \ldots. ``Primary hyperoxaluria should
  be suspected when urinary oxalate excretion exceeds 75 mg/day in
  adults without bowel dysfunction.''{[}5{]}

  \begin{itemize}
  \tightlist
  \item
    Type I
  \item
    Type II
  \item
    Type III (10\% of PH), Ox overproduction (not diet) UOx
    \textgreater{} 0.7 mmoL/1.73m\^{}2 /24 hours (BSA) Suble genetic
    differences characterize PH. Check:
  \item
    Ur GLYCOLATE
  \item
    L-glycerate
  \item
    4-hydro\ldots. (HOG)
  \item
    Di-hydro.. (DHG).
  \end{itemize}
\end{itemize}

\begin{enumerate}
\def\labelenumi{\arabic{enumi}.}
\setcounter{enumi}{1}
\tightlist
\item
  Enteric (\textasciitilde5\%, \textgreater80 mg/d, or 0.7-1.0 mmol/1.73
  m3/24 hours)) - GI problem: due to FAT MALABSORPTION and fat reaching
  colon (saponification of fatty acids with calcium) which eagerly binds
  with Ca in colon, leaving free Ox to be absorbed; symptom chronic
  diarrhea. Hard water (minerals) and soap (fatty acid) leaves residue
  hard to remove. This residue in COLON may also damage COLON MUCOSA,
  again allowing increased OX to be absorbed.
\end{enumerate}

ALSO, malabsorption fat (many possible causes): - bacteria \textbar{}
fat malabsorption \textbar{} chronic biliary or pancreas disorder
\textbar{} anything causing chronic diarrhea

Mechanism: Normal digestion, fats, bile acids absorbed in small
intestine. Leaving Ca, Mg and Ox to bind and be excreted. Anything goes
wrong with normal absorbtion and fat or bile or \ldots{} will continue
to colon. Ca will bind to these, less chance to bind to Ox. Also
possible colon tissue more porous to Ox, due to one of these
conditions.\\
Remark: castrated RATS have decreased UrOx; or with high dose
Finasteride!

\begin{enumerate}
\def\labelenumi{\arabic{enumi}.}
\setcounter{enumi}{2}
\item
  Diet (40-60 mg/d, 0.46-0.6mmol/1.73/d) - excess dietary Ox, low Ca
  consumption; lacking O. Formigenes; cell membrane Ox transport
  problem? {[}excess ascorbic acid \textgreater2 g/d increases Ox; 1
  orange 70-100mg ascorbic; 1 grapefruit, 88 mg; (100\% RDA = 60 mg?).
  Often temporary, but still can lead to chronic kidney disease
  (anti-freeze! or very high consumption of star fruit) ref:
  http://www.ncbi.nlm.nih.gov/pmc/articles/PMC4220346/ CASE STUDY:
  https://www.ncbi.nlm.nih.gov/pmc/articles/PMC5359106/
\item
  Idiopathic or Mild (40 - 60mg/d, 0.46-0.6 mmol/1.73/d) - spontaneous
  or unknown cause; might be diet or excess endogenous Ox by Liver; red
  blood cell transport; perhaps some genetic factor in one of these
  (``Differentiating between dietary and idiopathic hyperoxaluria can be
  difficult.'') Mayo: ``increased kidney excretion of oxalate in the
  absence of identifiable factors\ldots generally mild'' -- 13C can
  assist in intestinal absorption of Ox. -- If not genetic, Liver
  Biopsy..
\end{enumerate}

Hyperoxaluria - higher in summer (diet)

\begin{itemize}
\tightlist
\item
  ``Enteric hyperoxaluria is further subdivided into excess ingestion of
  large loads of oxalate-rich foods and/or increased small bowel wall
  permeability to oxalate uptake. \sout{Enteric hyperoxaluria leading to
  AON has numerous well-documented causes,} including consumption and
  juicing of oxalate-rich foods (spinach, parsley, peanuts, soy, and
  star fruit) in weight loss attempts, high intake of vitamin C
  supplementation, low vitamin B6 intake, toxic ingestions of ethylene
  glycol (antifreeze), and multiple causes of fat malabsorption (chronic
  pancreatitis, gastric bypass, medications inhibiting pancreatic
  lipase, and various causes of colitis including inflammatory bowel
  disease) {[}2, 5--9{]}.'' (source:
  https://www.ncbi.nlm.nih.gov/pmc/articles/PMC6323504/)
\end{itemize}

\textbf{HYPERURICOSURIA} - urinary uric acid excretion \textgreater800
mg/day (I did, pre allopuinol{[}2{]}) Cause: Too much diet purine, OR
endogenous uric acid production (vs

HYPERURICEMIA - excess uric acid in BLOOD. Increases risk b/c acts as
`seed' to CaOx xtals to form? (Claim: plasma oxalate s/d be ZERO;
\textgreater{} 30 is signal for dialysis) Normal: kidney should excrete
all uric acid; if load too high (meat protein), kidney may not be able
to eliminate quickly enough.

IDIOPATHIC X - have condition X, but w/o apparent underlying ETIOLOGY
(cause, origin, `reason') (ex: Idiopathic HyperOxaluria)

-- K --

KIDNEY ANATOMY (Khan video) Filters + Collects; maintains homeostasis
(status quo) inlcuding blood pH, BP, waste (urea), blood osmolarity --
Renal Cortex(shell), outer area of kidney, renal arttery splits into
nephrons where wrap round, dip into Medula. Renal Vein also collects
filtered blood from here. -- Renal Medula (middle), middle area where
neprhons dip down -- Renal Calyx, just below Medula, first place urine
appears -- Renal Pelvis, center most area where all Calyx combine, urine
flows to ureter -- Hilum, opening from kidney center where ureter, renal
vein, renal artery enter/leave See NEPHRON ANATOMY

\textbf{Kidney Biopsy} - invasive, useful to determine: acute tubular
injury, tubular damage, atrophy, interstitual mononuclear cell
inefficiency (?), oxalate crystal deposition in tubules/ and/or
interstitiute REF: KIReports

KETONE DIET. Claim: Switch from glucose (insulin) fuel to fat breakdown
for fuel (liver).\\
- Diet is minimal carbs, enough protein only, 70\% calories from
``healthy'' fats. These fats, also known as medium-chain Triglceries, do
NOT need bile to break down; good for pancreas insufficiency. - Liver
(lipase) will brake down triglerides to extract glycerol and convert to
fuel for muscles (glycogen). Advantages: ketone fuel (water soluable) do
not need protein to escort them around blood stream; easily pass
membranes to mitocondria. Less need for insulin at all. Insulin (and
related processes) associated with telling body to grow. If body thinks
less food around, it switches to preservation mode - keep every existing
cell alive, etc. Cancer can only use glucose, so ketone diet will
`starve' these cells. - Takes some effort to get Liver to transistion to
fat breakdown (triglericides). Starvation will do it, but not
pleasant.\\
- NOTE: vigorous exercise actually temporarily raises glucose! Walk
around post-exercise will mop this up.

-- L -- LAB TESTS (BLOOD) * BUN, 15-25 mg/dl; if too high, kidneys not
removing enough N from blood * Serum Creatinine, 0.5-1.3 mg/dl (not
considered accurate kidney fct test), can rise quickly (2.0, 3.0 \ldots)
but not as meaningful as GFR * However, change in SerCreatinine may
indicate change in GFR * Creatinine Clearance (24-hour urine) 90-130
ml/min (volume fo BLOOD that is cleared of Creatinine per min) * GFR
(preferred) - 80-120 ml/min/1.73 m\^{}2 (calculated) (preferred to Blood
Creatinine and does not require 24-hour urine); Stage 3 CKD less than 60
(lost ½ kidney function) LACTOSE - composed of 2 simple sugars
(galactose and glucose). To separate hydrolyze or in small intestine
exists lactase enyyzme (unless lactose-intoleratant?) . Commercial
lactose-free products: tablet with contain enzyme (ex:
Beta-galactosidase) to be taken BEFORE milk. Or, if lactose-free food,
then done at factory via different process. Claim: milk (12 g Lactose)
but yogurt (9 g) can't be. Ice cream 3g. I don't believe. LARGE
INTESTINE (or COLON + small pieces at beginning/end: CECUM and RECTUM) -
1.5 m long, food spends \textasciitilde16 hours here, * PROXIMAL - first
⅔ of COLON * DISTAL - last ⅓ of COLON * Cecum is beginning of COLON,
rich in bacteria. (Appendex attaches here) Over 99\% of the bacteria in
the gut flora are anaerobes,{[}4{]}{[}5{]}{[}6{]}{[}7{]}{[}8{]} but in
the cecum, aerobic bacteria reach high densities.{[}4{]}

``LEAKY GUT'' (discredited) - increased permeability of intenstine
(beyond normal). Due to illness? toxins? In this theory, responsible for
many illnesses (autism).

LEPTIN -Exercise decreases this chemical. A protein (`sensor', STAT3) is
sensitive to low leptin levels, triggering brain to `eat', via dopamine.
LEPTIN considered a regulator of energy storage.

LIPOPROTEIN Transport lipids, which are water insoluble. HUGE. Contain
LDLs, HDLs, Cholesterol etc. However, higher fat contents, less dense is
lipoprotein (less dense cream rises to top of milk): Per Am Heart, keep
Trig \textless100 (\textless150 is low risk)

LIVER BIOPSY: 1. Liver biopsy documenting alanine-glyoxylate
aminotransferase (AGT) activity below the normal reference range
confirming PH type 1 OR, 2. Liver biopsy documenting glyoxylate
reductase/hydroxypyruvate reductase (GR/HPR) activity below the normal
reference range confirming PH type 2

LOW-OXALATE DIET (LOX) - Variously mentioned as \textless40 mg,
\textless50 mg, \textless75 mg per day. LUMEN - open space in tubes,
where linings are often complex, multi-structure. INTRALUMIN

LYMPHATIC SYSTEM (vs.~Vascular) - Clear liquid between cells. Absorbed
in small intestine, transmitted to veins in neck (bypass LIVER) and
distributed in blood. Absorbs fat from intestine.
https://en.wikipedia.org/wiki/Lymphatic\_system\#Fat\_absorption

-- M -- Magnesium (Mg++) Role in stone prevention is DISUPUTED; Normal
urine, \textgreater= 70 mg/d of Mg. here is one paper:
https://www.ncbi.nlm.nih.gov/pubmed/16953377 Mg+2 - good at reducing
Ox-2 absorption, but we need less of it then Ca, so Ca is usually first
suggested. Because Mg correlates to Ox in most foods, Mg supplements are
recommended: easily digested. MgCitrate or supplemtns Others say no
clinical trials to support Mg supplements decrease risk in otherwise
healthy individuals. HOWEVER, people with diarera, bowel disease, fat
malabsorption may need Mg to replace loss;( not for any stone
prevention) mEq (milliquivalents) - arcane measure - millimoles x
VALANCE:\\
10 mEq Potassium = 1080 mg medullary nephrocalcinosis
https://radiopaedia.org/articles/medullary-nephrocalcinosis
https://www.zocdoc.com/answers/17936/what-is-medullary-nephrocalcinosis-how-can-it-be-treated
MET- Medical Expulsive Therapy - Use alpha-selective \textbar{}
alpha-blocker (ex: Tamsulosin) - relax ureter muscles (also relieves
some BHP symptom ?) - MRI measure kidney size/cysts or stones ?? (vs CT
and all the other ways?) NEPHRON ANATOMY- (khan video) -- basic
plumbing/cleaning element of kidney (filter + collect), containing
several tubing sections, including RENAL TUBULE where secretion and
reabsorption occur; \textasciitilde0.8-1.5 MM nephrons per kidney; blood
in; clean blood out; (100L/d/kidney) rest passes on as urine; -- blood
from renal artery flows into many arteriole arteries. Each leads to
glomerulus, where it winds itself into a strawberry-sized ball, and then
continues as efferent artery. At base of glomerulus is Bowman's capsule
or space where fluids collect (glucose, ions, white, red blood cells,
water, amino acids) but not large protein molecules where are too large.
Bowman's capsule encapsulates glomerulus. -- The interface between
glomerular and Bowen is at least 3 layers which limit what can and
can/not pass through. -- Diameters of arteriole and efferent arteries
affects amount of filtration. Ideally arteriole wide, to let a lot of
blood in, but efferent narrow forcing blood to slowly move through
glomerulus and filter completely. -- LOTS can go wrong! -- proximal
convoluted tubule, long winding tubing leaving the glomerulus with
water, Na, amino acids etc. from Bowen and reabsorbs \%65 of fluids;
then leaves renal cortex dives into Medulla for Loop of Henle (video)

-- RENAL TUBULAR ACIDOSIS (RTA) - failure of kidney to remove blood
acid, or add enough alkali to blood, so both blood and body accumulate
acid; also reduces UrCit) Since acid is NOT being eliminated, as normal,
the urine pH rises. When the kidney can no longer lower urine pH below
5.3 the situation is RTA.\\
-- Can live on total 300,000 nephrons (10\%) Claim: Kidney can produce
CITRATE, but if blood acidifies it will steal kidney CITRATE, leaving
less for Urinary CITRATE -- NEPHROTOXIN -- NSAIDS, COXIBS or heavy
metals, or rhubarb! (details) Measured by creatinine clearance (blood
test). NEPHROCALCINOSIS - calcium deposits in renal parenchymal, several
possible causes including hyperoxaluria (details)

NEPHROPATHY - kidney disease; do I have ``secondary oxalate
nephropathy''? Due to oxalate deposition. - \textbf{AON: Acute oxalate
nephropathy,} defined as acute kidney injury (AKI) due to biopsy-proven
calcium oxalate deposition, Prognosis: Poor. ie AON can cause AKI. -
\textbf{AKI: Acute kidney injury (aka acute kidney failure}
(https://en.wikipedia.org/wiki/Acute\_kidney\_injury) rapid failure of
kidney function (\textasciitilde{} 7 days) -

NEPHROLITHIASIS - stone in kidney (vs ureter).

NEPHROTIC SYNDROME -- high protein in urine, low serum protein; not same
thing as kidney failure. one symptom: foaming pee. see
https://en.wikipedia.org/wiki/Nephrotic\_syndrome

\begin{itemize}
\tightlist
\item
  O -
\end{itemize}

OSMOLES: Once dissolved, \# moles/L. Example 1 mole NaCL becomes 2
osmoles in water (both Na \& Cl). But 1 mole of glucose remains 1 osmole
because it does not dissociate.

OXALATE (Ox--) forms - Salt of Oxalic Acid, Ex: CaOx, MgOx; Also Ester
of Oxalic Acid. Oxalic acid is very reactive H2C2O4, removes rust.
Plants, no kidneys, use oxalate to store excess calcium from soil, water
etc as CaOx , esp.~In older leaves, store as CaOx (salt), less bioavail?
Younger leaves hold the Oxalic acid form, which has higher
bioavailability in humans. Oxalate 88 mg/1 mmoL 44 mg/f1 mEq. - OXALATE
Absorption - Theory? Absorption to blood depends on Ox solubility of
foods. Foods high in Ox may not have same bioavailability. Once Ox
absorbed \textbf{90\% of it will appear in urine 24-36 hours later.} -
OXALIC ACID (C2O4H2) dissociates into OX, conjugation base( C2O4 -2) and
H+. (pKc=1.27) or 2 H+ (pKa=4.27) Relatively strong acid (?); simplest
diCarboxylic acid (wiki: https://en.wikipedia.org/wiki/Carboxylic\_acid)
Known for transporting metals. pKa for -2 to 10 (weak acids) strong
acids pka less than -2 Case study:
http://www.revistanefrologia.com/en-publicacion-nefrologia-articulo-acute-renal-failure-due-oxalate-crystal-deposition-enteric-hyperoxaluria-X2013251411051070

\begin{itemize}
\tightlist
\item
  Oxalate measurement in urine: several
  labs:http://ndt.oxfordjournals.org/content/early/2011/03/30/ndt.gfr147.full.pdf
  Ex: Oxalate Kit, Trinity Biotech, Berkeley, NJ; 591-D, 100 Assays;
  need a spectrometer!\\
\item
  Ox measure in stool: Ixion Biotech Inc (Alac hug, FL) more exact then
  PCR Ox test? (Didn't women at Allex.. say very difficult?).
\item
  OXALATE catabolism - breakdown to smaller molecules
\item
  Claim: Stone-formers either (1) produce more Ox, or (2) absorb more
  food Ox, or (3) higher Ox interaction with food? Active (small
  intestine) vs.~Passive (gastro) Transport ???
\end{itemize}

``Sample'' high Ox foods: Almonds 31\% soluable Black Beans 5\% Cinnamon
(1789 mg/f100 g) 6\% Turmeric (1969 mg/100g) 91\%

\begin{itemize}
\tightlist
\item
  Why no Mg, K, Na oxalate stones? Because these are very soluble. KOx,
  NaOx very soluble); CaOx is not soluable (max is 5.0 - 6.7 g/L).
  However, in stomach, pH=2, CaOx is more soluble, so pass to
  intestines. Mg(AN=12) is 567x more soluble.\\
  Ex: Milk (Ca) + Rubarb (v. high Ox) → gritty, CaOx xtals in mouth.
  Ca-Ox in intestine=good (goes to feces), Ca-Ox in urine (bad)
\item
  FREE Oxalate is absorbed in LARGE COLON.
\item
  ``transient flood'' Even with adequate Calcium, 8\% of Ox is absorbed?
  (CLARIFY!) Of this, half is aborbed quickly, 1-2 hours, in small
  intestine. So if a meal contains 200 mg Ox (very high), expect 8 mg Ox
  (0.5 x 8\%) will absorbed in first 1-2 hours. This is too high for
  kidney. (source: Coe, quoting Holmes?) . However, 200 mg Ox spread
  throughout the day is much less harmful.
\end{itemize}

\textbf{OXALOSIS} - deposition of CaOx in kidney tissue
(vs.~nephrolithiasis CaOx xtals in kidney), in extreme cases can deposit
in eyes, bones, heart. (Source: Goldfarb, in Ox \& Hyperoxalia mtg,
think this is PH1)

OXALATE HYPERABSORPTION - explain \textbar{} where \textbar{} cause. One
of several models: (1) excess Ox absorption (2) Excess Ox production (3)
Excess retention of pre-cursor, but common xtals (4) Limited ability to
inhibit precursor xtal, etc etc

OXALOBACTER FORMIGENES - Bacteria loves Oxalate, will reduce urinary Ox.
Claim: \textasciitilde100\% of children have O.F. but only 60-80\% of
adults do (b/c antibiotics?){[}3{]} Paradox: most Ox absorbed in small
intestine (\textasciitilde2-6 hours), though can be absorbed in stomach
or colon (enteric), but OF is present only in colon! How does it
intercept Ox? Stool test? PCR?

OMBS (Oxalate Metabolizing Bacteria Series), other bacteria? Note: too
much Ox may actually inhibit O. Formigenes. Microbes in Urinary Tract
(??) may form ``adaptive extra-renal pathway'' to aid kidney eliminate
Ox??

OXALATE BASELINE - Value? fast for 10-12 hours, eat HIGH Ox, monitor
urine next 6-24 hours (normal Ur Ox \textless{} 40 mg/24 hours)

-- P --

\textbf{PANCREAS} ( Human pancreatic elastase 1 (E1) aids digestion
(exocrine) by breaking down fat in GI, specifically glycerol from
triglicerides (main form of human fat); Elastase remains undegraded
during intestinal transit. Therefore, its concentration in feces
reflects exocrine pancreatic function.

Actute pancreatitis:\\
During an inflammation of the pancreas, E1 is released into the
\textbf{bloodstream}. Thus the quantification of pancreatic elastase 1
in \textbf{serum} allows diagnosis or exclusion of acute
pancreatitis.{[}12{]}

Note: Dr Wei Wang said Elastase known for giving FALSE results.
\textbf{My value = 29 ug/g (stool?), 14FEB2019} - 200 µg elastase/g
stool indicate normal exocrine pancreatic function - 100-200 µg
elastase/g stool suggest mild to moderate pancreatic insuffiency - 100
µg elastase/g stool indicate exocrine pancreatic insufficiency
(\textbf{EPI}).

Causes: - Diagnosis/exclusion of exocrine pancreatic insufficiency
caused by e.g. Chronic Pancreatitis, Cystic Fibrosis (error in CFTR gene
=\textgreater{} cloride flow across membranes (serious), Diabetes
Mellitus, Cholelithiasis (Gallstones), ``Failure to Thrive'', Pancreatic
Cancer, Papillary Stenosis, tumor blocking in ducts.

Stool Fat 72-hour test: preparations:

\begin{itemize}
\tightlist
\item
  https://www.nationaljewish.org/treatment-programs/tests-procedures/gastroenterology/72-hour-fecal-fat-test
\item
  https://ltd.aruplab.com/tests/pub/2002356 (refrigerate)
\item
  https://www.healthline.com/health/fecal-fat\#results
\item
  https://www.mayocliniclabs.com/test-catalog/Specimen/8310
\end{itemize}

EPI details: " The need for PERT in PEI without symptoms is a matter of
debate, and randomized clinical trials on this issue are lacking. " {[}
NIH: Dignose \& Treat: diagram{]} (
https://www.ncbi.nlm.nih.gov/pmc/articles/PMC3831207/ )
https://en.wikipedia.org/wiki/Pancreatic\_elastase
https://labtestsonline.org/tests/stool-elastase

\textbf{CHRONIC PANCREATITIS} - long-standing inflammation of the
pancreas that alters the organ's normal structure and functions.{[}3{]}
It can present as episodes of acute inflammation. Enzymes can even start
working inside pancreas =\textgreater{} harms cells. - symptoms: pain,
diebetes I, fat in stool, wt loss - proximate cause: ducts obstruct,
calicification, swelling \ldots{}

Blood tests; The endoscopic retrograde cholangiopancreatography
(\textbf{ERCP}) test or CT or NMR; less invasive:
\href{https://www.ncbi.nlm.nih.gov/pmc/articles/PMC3292642/}{explaination
\& pictures}

\textbf{Radiological \& tests} JR: use this outine, put in references to
these procedures.

\begin{verbatim}
- CT, can miss
\end{verbatim}

CT: Note: subtle abnormalities of early to moderate chronic pancreatitis
are beyond its resolution, and a normal finding on this study does not
rule out chronic pancreatitis.

\begin{verbatim}
endoscopic ultrasound (EUS, through mouth), biopsy, sound - 
\end{verbatim}

\textbf{EUS}: https://pancreasfoundation.org/endoscopic-ultrasound-eus/
use sound, no radiation, through mouth, can take biopsy, very good for
studying pancreas, great for finding cancer. Requires general
anesthesia.

\begin{verbatim}
MRCP, form of MRI (magnetic resonance imaging).; last 20 years
\end{verbatim}

\begin{itemize}
\item
  MRCP(magnetic resonance cholangiopancreatography).\\
  now preferred to ERCP, non-invasive; last 20 years; shows flow in
  ducts. Done: winter 2019-2020. Repeat in 1 year.
  \href{https://en.wikipedia.org/wiki/Magnetic_resonance_cholangiopancreatography}{MRCP}

  ERCP - invasive, but can fix things, examine ducts, spot cancer.
\item
  \textbf{ERCP} (via throat, endoscopic retrograde
  cholangiopancreatography), very invasive, much more advanced procedure
  for GI doctors (often extra 1 year), now less done. Remove stones, put
  in stents, fix things.
  https://www.niddk.nih.gov/health-information/diagnostic-tests/endoscopic-retrograde-cholangiopancreatography
\item
  https://en.wikipedia.org/wiki/Endoscopic\_retrograde\_cholangiopancreatography
\end{itemize}

\textbf{References:}

Panceatisis confirmation:
(https://emedicine.medscape.com/article/181554-workup): \textbf{EPI}:
https://pancreasfoundation.org/patient-information/ailments-pancreas/exocrine-pancreatic-insufficiency-epi/

\begin{itemize}
\tightlist
\item
  https://en.wikipedia.org/wiki/Chronic\_pancreatitis
\item
  https://emedicine.medscape.com/article/181554-overview\#a5
\end{itemize}

\textbf{PATHOPHYSIOLOGY} - study of \ldots{} disease and how it
functions in body, i.e.~study of abnormal systems

PDI (Percussion Diuresis Inversion) Claim: Especially after SWL good way
to eliminate `sand' ``Hang upsidedown'' Inclinded plane on back (like
situps) or use chair, 10-15', back on floor, feet raised on chair. Hand
over head. Caution may raise BP PHYSIOLOGICALLY - how organ works, its
function, chemical and physical processes

POTASSIUM - Hypokalemia (K+ too low); Hyperkalemia (too high and needs
to be monitored) -- K+ concentration higher INSIDE Cell wall; Na+
concentration higher OUTSIDE cell wall. -- An increase in SERUM
Potassium ⇒ decreased kidney function. (too high can trigger heart
attack)

PROBIOTIC - contains LIVE bacteria

PREBIOTIC - nutrients for bacteria (but NO live) - garlic, onion,
beta-carotene, lots of foods do. CLAIM: oligosaccharides (oligo = few),
short chain simple sugars, only found in human (not cow) breast milk
feed gut bacteria, not the baby. Baby's breath contains Hydrogen, an
indicator than bacteria are eating (in colon). Foods absorbed in short
intestine, ie no bacteria, produce no hydrogen to exhale.

PHOSPHATE - kidneys should eliminate

\textbf{PROTEINURIA} - protein in urine; several possible causes,
including compromised kidneys, diabetes \ldots{} Normal kidneys readily
re-absorb blood protein and do not pass it to urine MONITOR urine
protein is key to manage GFR (?)

\begin{itemize}
\item
  Albuminuria - common protein found in urine when kidneys fail to
  filter; a large molecule should not be able to pass. No such thing as
  ``micro''

  \begin{itemize}
  \item
    A1 represents normal to mildly increased urinary albumin/creatinine
    ratio (\textless30 mg/g or \textless{} 3 mg/mmmol);
  \item
    A2 represents moderately increased urinary albumin/creatinine ratio
    (30--300 mg/g or 3--30 mg/mmmol, previously known
    asmicroalbuminuria); and
  \item
    A3 reflects severely increased urinary albumin/creatinine ratio
    \textgreater300 mg/g or \textgreater{} 30 mg/mmol).{[}1{]}
  \end{itemize}
\end{itemize}

\textbf{PURINES} - Considered a food source, associated with protein,
esp animal, that leads to higher serum uric acid (gout risk,
i.e.~crystaliziation) and urinary uric acid. Unknowns: role of dairy in
reducing? Plants have less, not zero purines. Meat \textbar{} Fish
\textbar{} eggs - esp.~Sardines, anchovy. TRUE? increase UrCa, decrease
Cit (both bad)

PHYTATE - Found in plants as tight bonds to store phosphorous (bran,
beans, nuts, cereals, seeds), varies with growing temp, fertilizer,
grain; Salt form of PHYTIC Acid. Urinary phytate reduces CaOx risk by
binding to Ca.

In PLANTS, inhibits growth in seeds grain legume till just right T and
moisture; blocks Phosphate absorption; binds with Ca, Mg, Fe, Zn; CLAIM:
people, esp developing countries where whole grain is dominate food, can
have depletion of minerals.

CLAIM: soaking, sprouting, fermenting will reduce phytic acid and
therefore increase digestion in human.

\begin{itemize}
\tightlist
\item
  Humans have limited enzyme phytase, ie can not digest, i.e.~so it is
  removed via feces. Considered ``Anti-nutrient.'' Found in high fiber
  diets. Claim: vegetarians have gut bacteria can break down phytate. In
  farm animal fed grains, this leads to excess soil Phosphorus. Ruminant
  animals can digest.\\
\item
  To increase bioavail and absorption of PHYTATE in intestine can
  cook\textbar acid\textbar ferment\textbar sprout (Role of bacteria?)
\item
  In acid form, PHYTATE BINDS to Ca, Fe, Zinc and others, decreasing
  human absorption of these metals, and can lead to deficiencies (ex:
  anemia) (Ca+2 binding is pH dependent) Some very high fiber African
  diets (mostly roots?) can yield excess phytate and nutritional
  problems. . Excess phytic acid ALSO inhibits enzymes (pepsin to digest
  protein, amylase to break starch to sugar, trypsin to digest protein).
\item
  But PHYTATE also correlates to Ox in many foods. In test-tube, PHYTATE
  inhibits Ca salt crystals. Only 5\% of PHYTATE reaches urine
  (``relatively non-absorbable). But, once in urine, PHYTATE considered
  good, since will bind to Ca in urine\\
\item
  PHYTATE or PHYTIC acid also plays role in intercellular; but source is
  NOT from diet and must be manufactured within cells.
\item
  Unclear! For non-stone formers, what is PHYTATE benefit, even if they
  can break it down? (into what?) For stone-formers sounds like good AND
  bad. In diet it correlates with Ox (BAD), but if it reaches urine then
  PHYTATE is good, binds Ca.
\item
  Since PHYTATE corrolaes with OX, should I eat or not? Some foods
  (soy,whole gains, beans, rice, cereals ) high in both OX and PHYTATE.
  Does PHYTATE mitigate high OX?
\item
  Benefits: (confused? anti-nutrient?) But PHYTATE itself can also lead
  to reduced blood glucose, decreased triglycerides, cholesterol, acts
  as antioxidant, remove \emph{excess} metals in body (Fe) and minerals
  related to kidney stones! (i.e.~provides `chelating') (GOOD). In
  contaminated soil, phytate will remove harmful metals!
\end{itemize}

PHYTASE - an enzyme with can break PHYTATE bonds, found in inactive form
in plants with PHYTATE. To activate, soak\textbar{} sprout \textbar{}
ferment. Humans with `healthy gut bacteria' produce some phytase enzyme
and possible to obtain nutrition from unsoaked grain. POTENTIAL RENAL
ACID LOAD (PRAL) - calc
(http://mathsurvey.org/potential-renal-acid-load-calculation.html) if
food will lead to acid or base PURINES - source of urinary uric acid.
30\% (estimate) of purines due to DIET{[}4{]}

\textbf{PTH} - Protein in blood from parathyroid gland, regulates serum
Ca and serum Phosphorous. \textbf{Hyperparathyroidism} is excess of PTH,
result of several possible factors including Vitamin D, presence of
disease etc. Normal: tight balance, if serum mineral is low, the gland
releases PTH which directs gut to absorb more mineral from food, or bone
to leach more minerals.(done in part with kidney also increasing
calcitriol vitamin D production.). Quickly readjusts to any deviations.

Abormal: excessive PTH even if no mineral deficiency?

\begin{itemize}
\tightlist
\item
  can be primary (gland producing excess PTH, uneeded, can be
\item
  fixed via surgery) or secondary (more complicated).
\item
  elevated PTH leads to increased urine Phosphate.
\item
  expect higher PTH would lead to \textbf{decreased} urCa, not the case
  in primary because urCa is effected by serum Ca conc (??). The net is
  increased ur Ca.
\item
  If urine has elevated phosphate and Calcium, stones are soon to follow
  and usually how hyperthyroidism is found.
\item
  Secondary, can be low calcium, low vitamin D in diet or other.
\item
  CKD, Kidney no longer able to produce ample calcitriol (PTH continues
  high).
\end{itemize}

Ref: (readable) (https://cjasn.asnjournals.org/content/10/7/1257)

{[}UT medical{]}
(http://www.utmedicalcenter.org/your-health/encyclopedia/test/003690/)

\begin{verbatim}
-   [hyperparathryoid]https://en.wikipedia.org/wiki/Hyperparathyroidism
-   [secondary](https://en.wikipedia.org/wiki/Secondary_hyperparathyroidism)
[medline:] ( http://bit.ly/2uwpDMq  )
\end{verbatim}

\begin{itemize}
\item
  May 2019 (Litholink)

  \begin{itemize}
  \item
    Serum P 3.7 (2.5-4 mg/dL)
  \item
    Serum Ca 9.4 (8.8 -10.2 mg/dL)
  \item
    UrCalcium 296 (less than 250 mg/day) and rising!
  \item
    UrP 1.222 (0.6 - 1.2 g/d)
  \item
    In diagnosed primary PTHT: elevated PTH. elevated serum
  \item
    Ca conc. lower serum Ph conc. wiki: test: ``intact'' PTH
  \item
    (excludes pieces of the molecule) normal (10 to 55 pg/mL)
  \item
    (July 2017) PTH, Intact 80 pg/mL (15 - 75
  \item
    pg/mL) H 2019: PTH 84 pg/mL high value could
  \item
    indicate Vit D problem or excess phosphorous (I am
  \item
    normal), due to chronic kidney disease. (July 2017, Dr
  \item
    Vanek, OHSU says to increase Vitamin D to 4,000 IU per
  \item
    day!) Also many factors can raise PTH
  \end{itemize}

  \textbf{2nd Hyperparathroid - Symptoms}

  \begin{itemize}
  \tightlist
  \item
    low serum Ca
  \item
    H serum Phosphate
  \item
    H ALP
  \item
    H PTH
  \end{itemize}

  \textbf{Tertiary Hyperparathryoid} chronic secondary hyperparathroid
  and gland no longer responds to changes in serum calcium. (bad news?

  \textbf{Parathryoid Adenoma} beign, almost always results in
  hyperparathryoid. Confirm by scan. Glad shouuld be in neck, but can
  migrate(?); so locating it is important. Surgery (only treatment) 95\%
  success.
\end{itemize}

\textbf{SALT} of an acid - add a metal Ex: HCL, salts: NaCL, KCl, CaCl
Other SALT examples: acid + base; metal + acid (Ca + oxalic acid),
others

\textbf{SLC26}- family of proteins, transport in kidney; exchange Ox for
sulfate\textbar chloride

\textbf{SODIUM} - for kidney to remove excess Na, must also excrete
Ca.\\
1 tsp SALT = 2325 mg SODIUM --- ``Unfortunately, faced with a chronic
excess of sodium to deal with, the kidneys can get worn out; sodium
levels in the blood then rise along with water needed to dilute it,
resulting in increased pressure on blood vessels and excess fluid
surrounding body tissues (read, swelling \ldots.'' (from Coe)

SUPERSATURATION - Unstable equilibrium in which concentrates exceeds
solubility but no precipitation occurs. Key = crystallization is
physical property and urine must stay below SS level (Coe, basis of
Litholink's SS). AUA Guideline Statement 6: ``Urinary supersaturation of
stone-forming salts is provided as part of the 24-hour urine analysis
panel from many commercially available laboratories, or it can be
calculated using a computer program. Supersaturation can guide and
monitor effectiveness of treatment and as such, is useful as part of the
initial 24-hour urine evaluation.46,47,55 ``.

RANDALL's PLAQUE - calcification (CaPhosphate) of subepithelial.

NEAR RENAL PAPILLA (termination of RENAL TUBUAL at connection to MINOR
CALYX, i.e base of Mendulla, aka Renal Pyramid. Less than 2mm thick.
Pre-cursor to stones? KUB can NOT distinguish between a stone and
RANDALL's PLAQUE (paper) or paper) use ureteroscopy. 1-2 mm, calcium
deposits; aka interstitial deposits (bet cells) of calcium. Is it CaP or
CaOx??

RENIN - If low Na (sweat, etc), kidney releases RENIN which triggers
LIVER to release Ang I, which triggers LUNGS to release Ang II which
tirrigers ADRENAL GLANDS (atop the kidney) Aldosterone\ldots. to tell
Kidney to retain Na and water (increasing BP), half-life 20'. ACE
reduces secretion of Aldosterone.

-- S --

SMALL INTESTINE (https://en.wikipedia.org/wiki/Small\_intestine) aka
small bowel - duodenum - jejunum, many things absorbed here. - ileum,
absorbs B12, bile acid, nutrients; next is large colon.

\textbf{STEATORRHEA} - https://en.wikipedia.org/wiki/Steatorrhea -
excess fat in feces (floats, smells, hard to flush, bulky); guideline
\textgreater{} 6 g/d stool fat

SUPERSATURATION (see Coe) - Concentration is above saturation level, but
no precipitation - unstable equilibrium. Meaning max concentration
before precipitation, ie max equilibrium -- just as many molecules
precipitate as those return to solution per unit time. Heat increases
solubility, for example sugar in water. Iif then cool SLOWLY , may
remain in solution. Shake it, crystals will form.

URINE, generally more stable (then sugar and water) and can remain SS
for days! BUT if making stones, we have SS urine and must lower the
concentration. Kidneys are water-efficient (ie reabsorb water) leaving
urine potentially SS w/ salts. Why? Fresh water once rare than now and
then much sodium in diet, so less Calcium in urine.

-- T --

Tamm-Horsfall protein (uromodulin) - low levels in urine may indicate
stones. Test?

-- U --

UUN24 - measure of urea in urine (normal: 7-20 g/d) see BUN.

UREA - Breakdown of proteins (nitrogen)

\textbf{URIC ACID} - in Urine, relatively insoluble, end product of
purine, in acid (less than 5.5) will precipitate (uric acid stones) BUT,
even in alkaline urine, excess uric acid can increase CaOx stone risk,
providing `seed' (process involves several steps, monoSodium Urate (MSU)
supersaturation. Also, MSU crystals in blood=GOUT).

-- Treatment for hyperuricemia : first raise urine pH with K+Cit- to
6.5-7.0, which may also (slowly) dissolve uric acid stones. If K+Cit-
does not reduce uric acid, then try Allopurinol.

-- Treatment for uric acid stones (uric acid lithiasis) - K+Cit-
especially at night to increase ``alkalide tide''
(https://en.wikipedia.org/wiki/Alkaline\_tide), blood pH rises for about
2 hours after meal. (So kidney increases UrCit ??) Stone formers have
decreased ``alkaline tide'' (?)

\textbf{URINE}- pH should be 5.5-7.0 Acidic urine will reduce citrate
(check?) the natural stone inhibitor. If sugar in blood, check for
Diabetes. To lower urine pH: grain, dairy, legumes, meat. To raise urine
pH: fruit \& vegetables. Urine is complex: thousands of organic
molecules.

\textbf{UREMIA} (``urine in the blood'') - originally referred to
execess serum urea (?) but now know that many other poisons in blood
when kidneys compromised.

\textbf{UROEPITHELIAL LINING} - Type of TRANSEPITHELIAL (``surrounds
lumin'') which lines LUMEN of urinary track (renal pelvis, ureter,
bladder, prostate) and is approx 5 cells deep. Can expand \textbar{}
contract. Barrier.

\textbf{Vitamin B6} - Rare, inherited disease in which body produces
excess oxalate, which binds with Ca in kidney tissue?? (oxalosis). In
some cases, B6 can help. Even in a few non-inherited cases, claim B6 may
work (note: `in a few') However, B6 supplementation has risks. Consult
doctor.

\textbf{Vitamin D} - Aids Ca absorption from intestine. When LOW, Ca
absorption \textasciitilde10-15\% of Ca food. When adequate, Ca
absorption 30-40\%. When Blood Calcium is LOW, PTH level rise and bones
leak Calcium. The PTH mechanism involves kidneys releasing activated
vitamin D, 1,25 dihydroxyvitamin D, or \textbf{calcitriol}.

Blood Test: 25(OH)D:\_\_\_ ng/ml x 2.5 = \_\_ nmol/L Institute of
Medicine 20 ng/ml MININUM - 50 ng/ml MAXIMUM Others say range: 40-60
ng/ml Claim: 2000 IU ---\textgreater30 ng/dL. 4000 IU ---\textgreater{}
50 ng/dL.

Low D? Symptoms include Bone Loss, Weaker Muscles, and (discovered last
30 years) all kinds of other bad things (few clinical trials)

Stone-formers: CLAIM: excess D will increase Ca absorption MORE THAN
non-stone formers. (Better check this\ldots. Not simple cause \& effect)
Leaving less Ca in gut to crystalize Ox. (So it too will be absorbed -
colon?) Avoid winter vacations where sun exposure abruptly changes from
nil to every day; every hour.

===== CT scan vs.~(readable)
http://www.controversies-and-consensus.com/lectures/moore.pdf

Video * Kahn: Kidney

Coe: http://kidneystones.uchicago.edu/fluid-prescription/ * Reread
(revised) * KIDNEY STONE TYPES * ANALYSE EVERY KIDNEY STONE * KIDNEY
STONE ANALYSIS - How Bad is It? * NEPHROCALCINOSIS * SUPERSATURATION *
CLINICAL SUPERSATURATION

\begin{center}\rule{0.5\linewidth}{0.5pt}\end{center}

{[}1{]} Mayo Clinic Communique, July/August 2013 {[}2{]} AUA, 2014
Guidelines \#16, ``allopurinol reduced the risk of recurrent calcium
oxalate stones in the setting of hyperuricosuria (urinary uric acid
excretion \textgreater800 mg/day) and normocalciuria.

{[}3{]}http://www.karger.com/Article/Abstract/80264 (Goldfarb - VA NY)
``There is evidence that the organism's absence, perhaps sometimes due
to courses of antibiotics, may be a cause of hyperoxaluria and stone
formation. In early investigations, patients not colonized with the
organism can be recolonized.'' {[}4{]}AUA, 2014, Guideline Statement 12:
``\ldots{} Nonetheless, if diet assessment suggests that purine intake
is contributory to high urinary uric acid, patients may benefit from
limiting high- and moderately-high purine containing foods. Although the
reported purine content of foods varies, ``high purine'' foods are
generally considered those containing more than 150 mg/3-ounce serving.
These include specific fish and seafood(anchovies, sardines, herring,
mackerel, scallops andmussels), water fowl, organ meats, glandular
tissue, gravies and meat extracts. ``Moderately-high'' sources of
purines include other shellfish and fish, game meats, mutton, beef,
pork, poultry and meat-based soups and broths.96,97 Note that an
individual may never or only rarely consume ``high purine'' foods but
may habitually consume large portions of foods in the ``moderately
high'' category.

5{]} Guideline Statement \#6 - 2014 AUA Guideline - Also, `` These
patients should be considered for referral for genetic testing and/or
specialized urine testing. REF\textless56\textgreater''

\hypertarget{misc-notes}{%
\subsubsection{Misc Notes}\label{misc-notes}}

\hypertarget{mon-07-dec-2020-notes-from-oxalate-hyperoxaluria-non-profit-group}{%
\subsubsection{Mon 07 Dec 2020, notes from Oxalate \& Hyperoxaluria
Non-Profit
group:}\label{mon-07-dec-2020-notes-from-oxalate-hyperoxaluria-non-profit-group}}

\begin{itemize}
\tightlist
\item
  water \& KCitrate is the standard.
\item
  PH1, some benefit from B6, others do not
\item
  FDA now considers UrOx a proxy for stones, stone burden.
\item
  PH1 kids are not in good shape; solutions include kidney transplant
  and sometimes BOTH kidney \& liver transplant.
\item
  \textbf{Lumasiran} new medicine (2020) promises some hope for PH1.
\end{itemize}
